% J2.3: limpiar es validar, marcar, tratar y registrar; nunca borrar en silencio.
\documentclass[tikz,border=4pt]{standalone}
\usepackage{resonancia}
\begin{document}
\begin{tikzpicture}
  \node[cajaAzul, minimum width=2.4cm] (v) at (0,0) {\textbf{1. Validar}\\[-1pt]{\footnotesize reglas de rango}};
  \node[cajaVerde, minimum width=2.4cm] (m) at (3.5,0) {\textbf{2. Marcar}\\[-1pt]{\footnotesize huecos, picos, planos}};
  \node[cajaNaranja, minimum width=2.4cm] (t) at (7.0,0) {\textbf{3. Tratar}\\[-1pt]{\footnotesize rellenar o excluir}};
  \node[caja, draw=rVioleta, fill=rVioleta!8, minimum width=2.4cm] (r) at (10.5,0) {\textbf{4. Registrar}\\[-1pt]{\footnotesize qué, cuántos y por qué}};
  \draw[flecha] (v) -- (m); \draw[flecha] (m) -- (t); \draw[flecha] (t) -- (r);
  \node[etiqueta, color=rGris, anchor=north] at (5.25,-0.95) {el dato crudo no se modifica: se añade una columna de marca y se trabaja sobre una copia};
\end{tikzpicture}
\end{document}
